%Vista lateral do robô
% (L4 desenhado inclinado para se ver o ângulo theta_C; a garra
%  é perpendicular a L4, por isso alpha = theta_C - 90 graus)
\begin{tikzpicture}[scale=0.3,
                    % Estilos
                        link/.style={line width=2.2pt, line cap=round},
                        junta/.style={circle, draw, fill=white, line width=1pt, inner sep=2.2pt},
                        cota/.style={gray!70!black, thin},
                        ang/.style={->, thin},
                        qang/.style={->, thin, red!75!black},
                        qcota/.style={red!60, thin, dashed}]

  % Coordenadas das juntas
    \coordinate (J1) at (0,0);
    \coordinate (J2) at (0,13.2);
    \coordinate (J3) at (3.8,23.5);
    \coordinate (J4) at (13.2,21.9);
    \coordinate (J5) at (20.06,24.40);
    \coordinate (J6) at (25.26,10.12);

  % Eixos
    \draw[->, thin] (-4,0) -- (31,0) node[below left] {$d$};
    \draw[->, thin] (J1) -- (0,29) node[left] {$h$};

  % Base
    \fill[gray!25] (-3,-0.8) rectangle (3,0);
    \draw[thick] (-5,0) -- (30.5,0);

  % Ângulos absolutos (theta_A, theta_B, theta_C e alpha, a partir da horizontal)
    \draw[cota, dashed] (J2) -- ++(6,0);
    \draw[ang] ($(J2)+(4,0)$) arc[start angle=0, end angle=69.75, radius=4] node[pos=0.55, right=1pt] {$\theta_A$};
    \draw[cota, dashed] (J3) -- ++(9,0);
    \draw[ang] ($(J3)+(8,0)$) arc[start angle=0, end angle=-9.66, radius=8] node[pos=0.5, right=2pt] {$\theta_B$};
    \draw[cota, dashed] (J4) -- ++(7,0);
    \draw[ang] ($(J4)+(4.5,0)$) arc[start angle=0, end angle=20, radius=4.5] node[pos=0.5, right=1pt] {$\theta_C$};
    \draw[cota, dashed] (J5) -- ++(5,0);
    \draw[ang] ($(J5)+(3.5,0)$) arc[start angle=0, end angle=-70, radius=3.5] node[pos=0.45, right=2pt] {$\alpha$};

  % Ângulos das juntas (q2, q3, q4, relativos ao elo anterior)
    \draw[qcota] (J2) -- ++(0,7.5);
    \draw[qang] ($(J2)+(0,6.5)$) arc[start angle=90, end angle=69.75, radius=6.5] node[pos=0.5, left=3pt] {$q_2 < 0$};
    \draw[qcota] (J3) -- ++(69.75:5);
    \draw[qang] ($(J3)+(69.75:4)$) arc[start angle=69.75, end angle=-9.66, radius=4] node[pos=0.35, above right=0pt] {$q_3 < 0$};
    \draw[qcota] (J4) -- ++(-9.66:5);
    \draw[qang] ($(J4)+(-9.66:3)$) arc[start angle=-9.66, end angle=20, radius=3];
    \node[red!75!black] at ($(J4)+(3.6,-3)$) {$q_4 > 0$};

  % Braços
    \draw[link] (J1)  -- node[pos=0.55, right] {$L_1$} (J2);
    \draw[link] (J2) -- node[right] {$L_2$} (J3);
    \draw[link] (J3) -- node[below left=0pt and -2pt] {$L_3$} (J4);
    \draw[link] (J4) -- node[above=1pt] {$L_4$} (J5);
    \draw[link] (J5) -- node[right=2pt] {$D = L_6 + L_7$} (J6);

  % Garra (perpendicular a L4)
    \begin{scope}[shift={(J6)}, rotate=20]
      \draw[line width=1.4pt] (-1.1,1.5) -- (-1.1,-0.5) (1.2,1.5) -- (1.2,-0.5);
      \draw[line width=1.4pt] (-1.1,1.5) -- (1.2,1.5);
    \end{scope}

  % Juntas
    \node[junta] at (J2) {};  \node[left=5pt]  at (J2) {J2 (ombro)};
    \node[junta] at (J3) {};  \node[above left=3pt and 0pt] at (J3) {J3 (cotovelo)};
    \node[junta] at (J4) {};  \node[below=5pt] at (J4) {J4 (pulso)};
    \node[junta] at (J5) {};  \node[above=4pt] at (J5) {J5};
    \fill (J6) circle (6pt);   \node[below right=4pt and 2pt] at (J6) {centro de preensão $(d, h)$};

  % Guias
    \draw[cota, dashed] (J6) -- (25.26,0);
    \draw[cota, dashed] (J6) -- (0,10.12);

\end{tikzpicture}